\documentclass[12pt]{article}
\usepackage[utf8]{inputenc}
\usepackage[T1]{fontenc}
\usepackage[margin=1in]{geometry}
\usepackage{setspace}
\usepackage{xcolor}
\usepackage[colorlinks=true, urlcolor=blue!60!black, linkcolor=blue!60!black]{hyperref}
\setlength{\parindent}{0em}
\setlength{\parskip}{0.6em}
\renewcommand{\baselinestretch}{1.0}
\begin{document}
\begin{center}
{\LARGE \textit{New York v. United States}, 505 U.S. 144 (1992)\par}
\vspace{0.5em}
{\large Case Brief --- Two-Page Edition\par}
\end{center}

\section{Memory Jogger}
Congress may encourage but not compel States to regulate: the anti-commandeering principle struck down the radioactive-waste Act's ``take title'' provision while upholding its spending and access incentives (\href{https://www.courtlistener.com/opinion/112768/new-york-v-united-states/}{505 U.S. 144}).

\section{Detailed Case Facts}
Facing a shortage of disposal sites for low-level radioactive waste, Congress enacted a state-brokered compromise in 1985. States that failed to provide for their waste faced monetary incentives, escalating denial of access to the three sited States' facilities, and, after 1996, a duty to take title to and possession of the waste and assume liability for damages. New York sought to site its own facility amid fierce local opposition.

\section{Procedural History}
New York and two counties sued the United States under the Tenth Amendment and Guarantee Clause. The District Court dismissed and the Second Circuit affirmed; the Supreme Court granted certiorari.

\section{Judicial Votes}
9--0 upholding the monetary and access incentives; 6--3 striking the take-title provision. O'Connor wrote for the Court; White (with Blackmun and Stevens) and Stevens concurred in part and dissented in part.

\section{Holding}
Congress may not commandeer state legislative processes by compelling States to enact and enforce a federal program. It may use conditional spending or conditional preemption, but it may not offer a choice between two coercive alternatives. The take-title provision is invalid and severable.

\section{Analysis of Opinions}
O'Connor reasoned that the Framers rejected the Articles' model of governing States in favor of regulating individuals directly, and that commandeering blurs political accountability. The incentives left States free to decline; take title ``crossed the line distinguishing encouragement from coercion.'' New York's earlier consent did not matter, since federalism protects liberty. White called the Act a state-negotiated bargain New York should be estopped from repudiating; Stevens rejected the anti-commandeering rule entirely.

\section{Examples: Future Applications}
\textbf{Same-Side:} barring legislatures from repealing sports-gambling bans is commandeering, as \textit{Murphy v. NCAA} held. \textbf{Opposite-Side:} letting EPA regulate private factories directly when a State declines to act is valid conditional preemption. \textbf{Fence-Sitter:} ``active efforts'' and record-keeping duties on state child-welfare agencies and courts, as in \textit{Haaland v. Brackeen}.

\section{Critique}
Critics note O'Connor's concession that the rule ``is not derived from the text of the Tenth Amendment itself,'' while functionalists doubt the accountability rationale and warn that the decision pushes Congress toward full preemption or coercive spending conditions.

\section{Key Quotations}
\begin{quote}
``In providing for a stronger central government, therefore, the Framers explicitly chose a Constitution that confers upon Congress the power to regulate individuals, not States.''
\end{quote}
\noindent --- Justice O'Connor, Opinion of the Court.

\end{document}
